\seccionreporte{Base de datos e historial de inferencias}{sec:base_de_datos}

\propositoseccion{%
  Cada inferencia debe persistirse para auditoría y consulta posterior. Documente esquema, índices y política de retención.%
}

\subsection{Motor y justificación}

Para este MVP, implementamos una arquitectura de persistencia dividida en dos componentes. Para la base de datos de los proyectos históricos (los embeddings), utilizamos \textbf{ChromaDB}. Elegimos este motor porque es una base de datos vectorial que funciona de manera muy eficiente en memoria y nos permite ejecutar búsquedas de similitud coseno sin requerir configuraciones complejas en el servidor.

Por otro lado, para registrar el historial de auditoría de inferencias, seleccionamos \textbf{SQLite}. La decisión de usar SQLite se basa en que, por el momento, nuestro único requerimiento es almacenar un log plano que el dashboard administrativo pueda consultar. Desplegar un clúster de PostgreSQL únicamente para esta tabla habría generado una sobrecarga innecesaria en el servidor. El sistema funciona perfectamente manteniendo los datos en un archivo local (\texttt{inference\_log.db}), con la proyección de migrar a PostgreSQL cuando la concurrencia de usuarios lo exija.

\subsection{Modelo de datos}

Esta es la estructura de la tabla principal que implementamos en SQLite:

\begin{tabular}{@{}llp{5.5cm}@{}}
  \toprule
  Campo & Tipo & Descripción \\
  \midrule
  \texttt{id} & INTEGER PRIMARY KEY & Identificador único autoincremental. \\
  \texttt{user\_id} & TEXT & ID del alumno o profesor (extraído del JWT). \\
  \texttt{timestamp} & REAL & Momento exacto de la inferencia (timestamp de Unix). \\
  \texttt{filename} & TEXT & Nombre original del PDF o archivo analizado. \\
  \texttt{score\_colision} & REAL & Porcentaje de plagio o similitud detectado. \\
  \texttt{nivel\_riesgo} & TEXT & Clasificación del riesgo ('Bajo', 'Medio' o 'Alto'). \\
  \texttt{academic\_alignment} & INTEGER & Calificación del 0 al 100 evaluando el cumplimiento del formato. \\
  \texttt{veredicto} & TEXT & Dictamen textual generado por el LLM. \\
  \bottomrule
\end{tabular}

\subsection{Operaciones}

\begin{itemize}
  \item \textbf{Insert:} Inyectamos un nuevo registro en la base de datos justo antes de retornar la respuesta HTTP al cliente en el endpoint \texttt{/pre-validate-proposal}.
  \item \textbf{Select paginado:} Implementamos la ruta \texttt{GET /inference-history?limit=15\&offset=0}, la cual permite al frontend consumir los registros de monitoreo de manera fluida y sin sobrecargar la red.
\end{itemize}

\subsection{Privacidad}

Debido a que las propuestas pueden contener nombres de empresas, asesores o alumnos, procesamos todo el texto utilizando el modelo \texttt{es\_core\_news\_sm} de SpaCy antes de almacenarlo o evaluarlo. SpaCy localiza las entidades nombradas (NER) correspondientes a personas o ubicaciones, y las elimina del texto antes de generar los embeddings o enviar el contexto a Ollama. De esta manera, aseguramos el cumplimiento de las normativas de privacidad y evitamos exponer información sensible.
